\cotefile{158}
\scanpage{36}\mspage{16}

\[
\bar\varphi^{h}=1,
\]
donc $\bar\varphi$ est inversible. Mais alors $\deleted{\unsure{\bar g}}$
$\bar g\overset{\mathrm{déf}}{=}\bar\varphi^\beta$ est inversible (et $\unsure{=} $ d'ordre fini),
\deleted{Or dans un \unsure{monoïde}} et d'autre part il est minimal, comme
image de $\deleted{\unsure{u}}\varphi^\beta$ qui est minimal. \deleted{Ainsi,}
Mais si dans un monoïde $\bar M$ \deleted{tel que} \inserted{où}
on a la loi $[\bar u\bar f=\bar v\bar f\Longrightarrow\bar u=\bar v]$, on a
un élément minimal et inversible \inserted{$\unsure{\sigma|\varphi}$}, on
a vu que $\bar M$ est un groupe. Car
$\deleted{\unsure{\exists}}\forall u\in\bar M$, $\exists v\in M$ tel que $\deleted{\unsure{uv}}$
$\sigma=uv$ i.e. $u\underbrace{(v\sigma^{-1})}_{w}=1.$, et alors
$p=wu$ satisfait $p^2=p$ \deleted{d'où} $p.p=1.p$
i.e. $p=1$, \deleted{\unsure{donc}} $w$ est un inverse de
$u$, \deleted{\unsure{une}} tout élément est inversible.

On a prouvé :

\uline{Proposition 3} Soit $M$ un monoïde ayant
un \deleted{\unsure{ajo|objet}} élément minimal $\psi$, et soit
$\varphi\in M$ tel qu'il existe $\beta\in\mathbb N$ avec $\varphi^\beta$ minimal
(par exemple $\varphi$ minimal \deleted{\unsure{i.e. Si}}, \unsure{p. ex|p ce|par ex.}
$\varphi=\psi$). S'il existe $r,s\in\mathbb N$ avec $r\ne s$, et
$\varphi^r\psi=\varphi^s\psi$,

% Printed form upside down relative to the mathematics.
% letterhead
{\footnotesize
\begin{tabular}{@{}l@{\quad}l@{}}
\deleted{MINISTERE DE L'EDUCATION NATIONALE} & \deleted{AVIGNON, le}\\
\deleted{ACADEMIE D'AIX-MARSEILLE} &\\
\hfil\rule{1.5cm}{0.4pt}\hfil &\\
\deleted{Inspection Académique} & \deleted{L'Inspecteur d'académie,}\\
\deleted{de Vaucluse} & \deleted{directeur des services départementaux}\\
\hfil\rule{1.5cm}{0.4pt}\hfil & \deleted{de l'Education Nationale de Vaucluse,}\\
\deleted{Tél. : 90 82 99 76} &\\
\deleted{Référence : Secrétariat n\textsuperscript{o}} & \deleted{à M}
\end{tabular}}
% end letterhead

% archive sheet number: 35

\scanpage{37}\mspage{15 bis}

\[
\deleted{=(\alpha w)^{-1}(\alpha g\unsure{u|v})(\alpha gv)^{-1}(\alpha gfu)
 ((\alpha w)^{-1}(\alpha gfu))^{-1}=1}
\]

\textcircled{0} $M$ pseudocofiltrant i.e. $\forall u,v\quad\exists u',v'$ avec $uu'=vv'$

\textcircled{1} $\forall u$, $\exists u'$ tel que $uu'\sim\psi$ i.e. $\exists p,q\in\mathbb N$ tels que
\[
\varphi^p uu'=\varphi^q\psi
\]

\textcircled{2} $\exists a\xrightarrow{\lambda}e=a\psi$ i.e. $\exists\lambda$ tel que $\psi\lambda\sim1$ i.e.
\[
\exists\alpha,\beta\in\mathbb N\text{ tels que }\quad\boxed{\varphi^\alpha\psi\lambda=\varphi^\beta}
\]

[ \uline{NB} Si $\psi$ est \uline{minimal}, alors \textcircled{0} et \textcircled{1} sont
satisfaits, et \textcircled{2} est satisfait sss \deleted{\unsure{Question 3}}
$\exists\beta\in\mathbb N$ tel que $\varphi^\beta$ soit minimal \unsure{\}.|?|(.}. Et
pour un tel $\beta$, $\exists\lambda\in M$ (et $\unsure{=} $ $\lambda$ \unsure{minimal|minim.?})
tel que $\psi\lambda=\varphi^\beta$ ]

\textcircled{3} $\forall r,s\in\mathbb N$ avec $\varphi^r\psi=\varphi^s\psi$, on a $r=s$.

[ NB Si $\psi$ est minimal, alors la relation
\[
\mathscr R_0:\qquad u\overset{R_0}{\sim}v\quad\Longleftrightarrow\quad u\psi=v\psi
\]
est aussi la relation d'équivalence définie par
l'application
\[
\begin{aligned}
M&\longrightarrow \unsure{\mathrm{Ep}|\mathrm{E}p}(M_0)\\
u&\longmapsto u_{M_0}=(\lambda\longmapsto u\lambda:M_0\longrightarrow M_0),
\end{aligned}
\]
(puisque $u\psi=v\psi\Longrightarrow u\psi\deleted{\unsure{\beta}}\alpha=v\psi\deleted{\unsure{\beta}}\alpha\quad\forall\alpha\in M$,
or $\psi M=M_0$ i.e. l'ens des $\psi\alpha$ est l'ens des
$\lambda\in M_0$). Soit $\bar M\simeq M/R_0\subset\mathrm{Ep}(M_0)$ le
monoïde image. \deleted{Ainsi, le} C'est un monoïde
dans lequel on a : $\bar u\bar f=\bar v\bar f\Longrightarrow\bar u=\bar v$, et
il est pseudocofiltrant, comme image de
$M$. (Il n'est filtrant que s'il est réduit à $1$.)
Ainsi la relation $\varphi^r\psi=\varphi^s\psi$ dans \textcircled{3} équivaut
à $\bar\varphi^r=\bar\varphi^s$, où $\bar\varphi$ est l'image de $\varphi$ dans $\bar M$.
Si on a $r=s+h$, $h\geq1$, \deleted{on en conclut} \inserted{cela équivaut} à

% archive sheet number: 36

\scanpage{38}\mspage{17}

alors l'image $\bar M$ de $M$ dans $\unsure{\mathrm{Ep}|\mathrm{E}p}(\deleted{M_0}E_0)$ ($\deleted{M_0}E_0=$ ens. des
éléments minimaux de $M$) est un groupe. Donc
si $\inserted{Y=}B_M$ est $1$-connexe, ce groupe est réduit à $\{1\}$,
ce qui signifie que pour $\deleted{\unsure{\psi}}\lambda\in\deleted{M_0}E_0$ (p. ex.
pour $\lambda=\psi$), \deleted{et $u$} et $u\in M$, on a $u\lambda=\lambda$,
ce qui implique que $M$ est cofiltrant.

Dans le cas contraire $(\varphi^r\psi=\varphi^s\psi\Longrightarrow r=s)$
prenant la relation d'équivalence \inserted{inv. à droite} $R_\varphi$ \inserted{sur $M$} définie par
$u\overset{R_\varphi}{\sim}v\Longleftrightarrow\exists p,q\in\mathbb N$ avec $\varphi^p u=\varphi^q v$, et
le quotient $F=M/R_\varphi$, \unsure{vu|considéré} comme $M$-ens.
à droite, i.e. comme objet de $Y^\wedge$, alors
$Y_{/F}$ est \inserted{(pseudocofiltrant et)} $0$-connexe (tous ses objets sont $\unsure{=}$
équivalents), mais \deleted{\unsure{il|elle}} n'est pas $1$-connexe ;
plus précisément son groupe fondamental
est isomorphe à $\pi_1(M_0)$, où \unsure{1}
\[
M_0=\{f\in M\mid\psi f\overset{R_\varphi}{\sim}\psi\},
\]
et admet un groupe quotient isomorphe à $\mathbb Z$,
grâce à l'hom. surjectif
\diagram{\begin{tikzcd}[ampersand replacement=\&]
\pi_1(M_0)=G_0\arrow[r,"\delta_0"] \& \mathbb Z
\end{tikzcd}}
donné sur $M_0$ par
\[
\delta(f)=q-p\qquad\text{si}\qquad\varphi^p\underline{\psi f}=\varphi^q\psi.
\]

[ \uline{NB} \deleted{On} \deleted{\unsure{Il exi}} Comme $\psi\in M_0$, on a $\varphi\psi\in M_0$,
\deleted{\unsure{de même}} donc $\exists f$ tel que $\psi f=\varphi\psi$,
et alors $f\in M_0$ et $\delta(f)=1$ ] Donc on a $\unsure{=}$
$H^{1}(X,\unsure{\mathbb Q|Q|\mathbb Z})\ne0$.

% archive sheet number: 37

\scanpage{39}
% letterhead
% The typescript reads with the scan turned through 180 degrees.
ACADEMIE D'AIX-MARSEILLE\\
\rule{2cm}{0.4pt}\\
INSPECTION ACADEMIQUE\\
DE VAUCLUSE\\
\rule{0.9cm}{0.4pt}\\

C.D.A.C.N\textsuperscript{o}16\\

AVIGNON le 6 décembre 1990\\

Jacques ROUX, responsable\\
départemental de l'Action\\
Culturelle,\\

à\\
M\\

\uline{OBJET} : Projets d'Action Educative Collège \unsure{.}\\
Année scolaire 1990/1991\\

L'attribution des aides aux P.A.E. se déroulera cette année en deux temps :

DECEMBRE \quad - instruction des dossiers par l'Inspection \unsure{Académique|A\.{c}\.{a}démique}

propositions à la M.I.A.D.A.C.

JANVIER \quad - Monsieur le Recteur attribue une aide globalisée à Monsieur l'Inspecteur d'Académie.

Une commission se réunit vers le 15 janvier 1991, afin d'assurer la ventilation des aides au niveau du département. Je vous demanderai alors de participer à cette Commission.

P/L'Inspecteur d'Académie\\
\handwritten{\unsure{Roux}}\\
% The next typed line is cut by the edge of the scan.
\illegible{1}

% archive sheet number: 38

\scanpage{40}\mspage{18}

\deleted{\unsure{Dans le|Donc}} \inserted{Ainsi}

\uline{Corollaire} Soit $M$ un monoïde ayant un
objet minimal. Si \deleted{$M$ n'est pas \unsure{co}} $Y=B_M$
est tel que \deleted{tous} les objets $0$-connexes de
$Y^\wedge$ sont $1$-connexes (ou \unsure{simplement|seulement}, que
leur $H^1(x,\unsure{\mathbb Q|Q})=0$), alors $M$ est cofiltrant.

\medskip\centerline{\rule{1.4cm}{0.4pt}}\medskip

\uline{Théorème} Soit $X$ \inserted{$\in\mathrm{Ob}\,\mathrm{Cat}$} une catégorie pseudo-cofiltrante
(hyp. vérifiée si les objets $0$-connexes \inserted{$F$} de $X^\wedge$
sont $1$-connexes, ou si seulement il existe
un groupe \inserted{abélien} $k\ne0$ tel que \unsure{pour|que|l'on} \deleted{\unsure{ait}} $H^1(F,k)\ne0$).
On suppose que $X$ admet \deleted{des \unsure{obj|objets}} un plus
\deleted{objet} petit objet $e$, et que $X_{/e}$ (qui est
également pseudo-cofiltrant) en admet un également,
ou ce qui revient au $=$, que le monoïde
$M=\mathrm{End}_X(e)$ admet un objet minimal $\varphi$. (NB
ces deux hypothèses de minimalité sont satisfaites
si $\mathrm{Ord}\,X$ et $M=\deleted{\unsure{\lambda|\mathrm{End}}}\mathrm{End}_X(e)$ sont finis, en particulier
si $X$ est finie.) Considérons l'opération
de $\mathbb N$ sur $e$ (identifié à un objet de $X^\wedge$)
définie par $\varphi$, $T_n$ opérant par $\varphi^n$, et distinguons
deux cas :

a) Cette opération \deleted{se} fait ``str. sans pts fixes'', i.e.
$\mathbb N\times e\longrightarrow e\times e$ défini par $(n,u)\longmapsto(\varphi^n u,u)$ est
mono. Il revient au $=$ de dire que pour

% archive sheet number: 39

\scanpage{41}
% letterhead
% The typescript reads with the scan turned through 180 degrees.
ACADEMIE D'AIX-MARSEILLE\\
\rule{2cm}{0.4pt}\\
INSPECTION ACADEMIQUE\\
DE VAUCLUSE\\
\rule{0.9cm}{0.4pt}\\

C.D.A.C.N\textsuperscript{o}16\\

AVIGNON le 6 décembre 1990\\

Jacques ROUX, responsable\\
départemental de l'Action\\
Culturelle,\\

à\\
M\\

\uline{OBJET} : Projets d'Action Educative Collège \unsure{.}\\
Année scolaire 1990/1991\\

L'attribution des aides aux P.A.E. se déroulera cette année en deux temps :

DECEMBRE \quad - instruction des dossiers par l'Inspection \unsure{Académique|A\.{c}\.{a}démique}

propositions à la M.I.A.D.A.C.

JANVIER \quad - Monsieur le Recteur attribue une aide globalisée à Monsieur l'Inspecteur d'Académie.

Une commission se réunit vers le 15 janvier 1991, afin d'assurer la ventilation des aides au niveau du département. Je vous demanderai alors de participer à cette Commission.

P/L'Inspecteur d'Académie\\
\handwritten{\unsure{Roux}}\\
% The next typed line is cut by the edge of the scan.
\illegible{1}

% archive sheet number: 40

\scanpage{42}\mspage{19}

$u\in M$ et pour $p,q\in\mathbb N$, $\varphi^pu=\varphi^qu\Longrightarrow p=q$,
ou encore \deleted{que $\varphi^p$} (se ramenant au cas $u=\varphi$
puisque $\deleted{\unsure{\exists}}u'$ avec $uu'=\varphi$) $\varphi^p=\varphi^q\Longrightarrow p=q$.
Dans ce cas, par des principes généraux,
\deleted{le quotient} l'objet $e$ de $X^\wedge$ étant $\unsure{0|2}$-connexe,
l'objet quotient $F=e/\mathbb N=e/\varphi$ est $0$-connexe,
et son $\pi_1$ admet un groupe quotient isomorphe
à $\mathbb Z$, donc $\deleted{\unsure{F}}X_{/F}$ n'est pas $1$-connexe,
et $=\ H^1(F,\deleted{\unsure{Q|\mathbb Z}}k)\ne0$ pour tt groupe
abélien $k\ne0$.

b) Il existe $p\ne q$ tel que $\varphi^p=\varphi^q$. Alors
\deleted{l'opération} l'image de $M$ dans l'ens.
$\unsure{\mathrm{Ep}|\mathrm{E}p}(E)$ des épimorphismes de $E=$ ens des
éléments minimaux de $E$ est un \uline{groupe}
$\Gamma$, et si $X_0$ désigne l'ouvert de $X$
formé des plus petits objets de $X$, alors
\[
\deleted{\unsure{\pi}}\quad \pi_1(X_0,e)\simeq\pi_1(M)\xrightarrow{\text{épi}}\Gamma.
\]
Donc si $X_0$ est $1$-connexe, alors $\Gamma=\{1\}$, donc
on a $u\varphi=\varphi\quad\forall u\in M$, et $M$, \deleted{et} donc
aussi $X_{\deleted{\unsure{0}}}$, \deleted{\unsure{et $X$}} sont cofiltrants. \deleted{\illegible{1}}
% slanted
\margin{En particulier, on a $\varphi^2=\varphi$ pour tt $\varphi$ minimal.}

\uline{Corollaire} Si \unsure{tout|$\mapsto$|tous les} objet de $X^\wedge$ est $1$-connexe,
alors $X$ est cofiltrante,

% archive sheet number: 41

\scanpage{43}
% letterhead
% The typescript reads with the scan turned through 180 degrees.
ACADEMIE D'AIX-MARSEILLE\\
\rule{2cm}{0.4pt}\\
INSPECTION ACADEMIQUE\\
DE VAUCLUSE\\
\rule{0.9cm}{0.4pt}\\

C.D.A.C.N\textsuperscript{o}16\\

AVIGNON le 6 décembre 1990\\

Jacques ROUX, responsable\\
départemental de l'Action\\
Culturelle,\\

à\\
M\\

\uline{OBJET} : Projets d'Action Educative Collège \unsure{.}\\
Année scolaire 1990/1991\\

L'attribution des aides aux P.A.E. se\unsure{.} déroulera cette année en deux temps :

DECEMBRE \quad - instruction des dossiers par l'Inspection \unsure{Académique|A\.{c}\.{a}démique}

propositions à la M.I.A.D.A.C.

JANVIER \quad - Monsieur le Recteur attribue une aide globalisée à Monsieur l'Inspecteur d'Académie.

Une commission se réunit vers le 15 janvier 1991, afin d'assurer la ventilation des aides au niveau du département. Je vous demanderai alors de participer à cette Commission.

P/L'Inspecteur d'Académie\\
\handwritten{\unsure{Roux}}\\
% The next typed line is cut by the edge of the scan.
\illegible{1}

% archive sheet number: 42

\scanpage{45}\mspage{1}

\uline{Proposition} \uline{1} Soit $M$ un monoïde fini, tel que

a) le groupe enveloppant de $M$ soit $\simeq\{1\}$, et

b) $M$ est pseudo-cofiltrant \deleted{\unsure{ou} pseudo-cofiltrant}.

Alors il existe un \deleted{unique} élément $p$ de $M$ tel
que l'on ait, pour tout $u\in M$
\[
up\deleted{\unsure{=pu|=up}}=p,\qquad(\text{en particulier. }p^2=p),\tag*{$(\unsure{\times|*})$}
\]
\deleted{en particulier} \inserted{donc} $M$ est \deleted{\unsure{filtrant et|filtrant}} cofiltrant.

\deleted{\uline{Dém} L'unicité de $p$ est claire, car si $p'$ \deleted{\unsure{est}} \inserted{$\in M$} \deleted{un}
\deleted{objet avec} \unsure{la $=$e propriété|les $=$s propriétés} que $p$, on aura
(faisant $u\deleted{\unsure{\varphi}}=p'$ dans $(\times)$)}
\[
\deleted{p'p=pp'=p}
\]
\deleted{et symétriquement $\deleted{p'p=pp'}\quad pp'=p'p=p'$, donc
$p=p'$.}

\deleted{On peut supposer} $M$ pseudocofiltrant \deleted{\illegible{1}}
implique que pour un nb fini d'\deleted{objets}\inserted{éléments}
$u_i$ de $M$, il existe un \deleted{objet}\inserted{élément} $p$ de $M$ tel
\deleted{quelque soit} \unsure{que} pour tout $i$, $p$ se factorise en
\[
p=u_i v_i.
\]
Comme $M$ est fini, il existe donc un $p$
tel que pour tout $u\in M$, on ait une
factorisation
\[
p=uv_{(u)},\tag*{$(\times\times)$}
\]
où $v_{(u)}$ est un objet de $M$. \deleted{Cette propriété} \inserted{L'existence}
d'un tel $p$ implique visiblement que $M$
est cofiltrant, donc équivaut (pour $M$ fini)

% archive sheet number: 44

\scanpage{46}\mspage{2}

% letterhead
{\footnotesize
\begin{tabular}{@{}l@{\quad}l@{}}
\deleted{MINISTERE DE L'EDUCATION NATIONALE} & \deleted{AVIGNON, le}\\
\deleted{ACADEMIE D'AIX-MARSEILLE} &\\
\hfil\rule{1.5cm}{0.4pt}\hfil &\\
\deleted{Inspection Académique} & \deleted{L'Inspecteur d'académie,}\\
\deleted{de Vaucluse} & \deleted{directeur des services départementaux}\\
\hfil\rule{1.5cm}{0.4pt}\hfil & \deleted{de l'Education Nationale de Vaucluse,}\\
\deleted{Tél. : 90 82 99 76} &\\
\deleted{Référence : Secrétariat n\textsuperscript{o}} & \deleted{à M}
\end{tabular}}
% end letterhead

à l'hypothèse b.

\deleted{OPS} $M$ réalisé comme monoïde d'endomorphismes
d'un ensemble \uline{fini} $E$ (p. ex.
par la représentation régulière gauche de $M$).
La relation $(\times\times)$ implique
\[
\mathrm{Im}\,p\subset\mathrm{Im}\,u\qquad\forall u\in M,
\]
donc \inserted{$E_0=$} $\mathrm{Im}\,p$ est le plus petit des $\mathrm{Im}\,u$,
pour $u\in M$. Soit $u\in M$, on a
\[
u(E_0)=up(E)=\mathrm{Im}\,up)\supset E_0,
\]
\margin{\unsure{?}}et comme $E_0$ est fini, on en conclut que
$u(E_0)=E_0$ et que $u\mid E_0$ est un autom. de
l'ens. $E_0$. Soit $u_{E_0}$ cet autom. L'application
\[
u\longmapsto u_{E_0}\ :\quad M\longrightarrow\mathrm{Aut}(E_0)
\]
est un hom de $M$ dans un groupe. L'hyp.
a) sur $M$ implique donc qu'on a $u_{E_0}=\mathrm{id}_{E_0}$
$\forall u\in M$. On en conclut aussi que
\[
E_0=\text{ens. des pts fixes de }M\text{ opérant dans } E,
\]
car si $x\in E\setminus E_0$, comme $p(x)\in E_0$, \deleted{on} a $p(x)\ne x$,

% archive sheet number: 45

\scanpage{47}\mspage{3}

% letterhead
{\footnotesize
\begin{tabular}{@{}l@{\quad}l@{}}
\deleted{MINISTERE DE L'EDUCATION NATIONALE} & \deleted{AVIGNON, le}\\
\deleted{ACADEMIE D'AIX-MARSEILLE} &\\
\hfil\rule{1.5cm}{0.4pt}\hfil &\\
\deleted{Inspection Académique} & \deleted{L'Inspecteur d'académie,}\\
\deleted{de Vaucluse} & \deleted{directeur des services départementaux}\\
\hfil\rule{1.5cm}{0.4pt}\hfil & \deleted{de l'Education Nationale de Vaucluse,}\\
\deleted{Tél. : 90 82 99 76} &\\
\deleted{Référence : Secrétariat n\textsuperscript{o}} & \deleted{à M}
\end{tabular}}
% end letterhead

donc $x$ n'est pas pt fixe pour $M$.

\deleted{On en conclut aussi que}
Le résultat $u_{E_0}=\mathrm{id}_{E_0}$ \inserted{$\forall u\in M$} s'écrit aussi
\[
\begin{aligned}
up(\unsure{x|n})&=p(x)\qquad\forall x\in E,\quad\text{i.e.}\\
up&=p\qquad\forall u\in M,
\end{aligned}
\]
ce qui prouve déjà que $M$ est cofiltrant
(car $u,v\in M$, on a $up=vp$, les deux mbres
étant égaux à $p$), \uline{qed}.

\deleted{J'ai envie de prouver ceci :}

\deleted{\uline{Proposition} 2 \inserted{(?)} Soit $X$ une catégorie \uline{1-connexe}
\uline{finie}, ayant un \deleted{objet} plus grand objet
$a$, et pseudo-cofiltrante. Alors $X$ est
cofiltrante.}

\deleted{On retrouve $(\unsure{\pm|+|\ddagger})$ la prop 1 en prenant
$X$ réduite à un seul objet. Je vais
essayer de calquer la dém. de prop 1.
Je note d'abord que, dans la catégorie
$X_{/a}$, \inserted{\unsure{quel que soit|quel que|qu'il est} \unsure{ds. cat.|cat.} $X$ finie,} il y a un plus petit objet,}

% archive sheet number: 46

\scanpage{48}\mspage{4}

\deleted{soit $(e_{\deleted{\unsure{a}}},p)$, avec}
\[
\deleted{p:e\longrightarrow a.}
\]

\deleted{Si $b$ est dans $X$, on sait \inserted{de $=$} que \inserted{dans} la
catégorie $X_{/b}$ il y a un plus petit
objet, et on voit \deleted{\unsure{tout de suite}} \inserted{facilement} qu'on
peut prendre sa source égale à $e$,
donc le prendre de la forme}
\[
\deleted{\deleted{\unsure{\longrightarrow b}}\qquad e\xrightarrow{p_b}b.}
\]

\deleted{Mais \deleted{\unsure{a pr}} il n'y a pas unicité \unsure{de|des}
$p_b$, et pour assurer une unicité, si
faire se peut, on voudrait mettre
en relation le choix de $p_b$, avec
celui (déjà fait) de $p_a=p$. Comme
$a$ est un plus grd objet de $X$,
il existe}
\[
\deleted{u:b\xrightarrow{\unsure{\circ}}a\,\unsure{x|r}},
\]
\deleted{et on voit facilement aussi que pour
$u$ fixé, on peut choisir $p_b$ de
la forme $p_{bu}$, avec}
\diagram{\begin{tikzcd}[ampersand replacement=\&,row sep=small,column sep=small,execute at end picture={\draw (current bounding box.south west)--(current bounding box.north east);}]
\& \deleted{b}\arrow[dr,"\deleted{u}"] \& {}\\
\deleted{e}\arrow[ur,"\deleted{p_u}"]\arrow[rr,"\deleted{p}"'] \& {} \& \deleted{a}
\end{tikzcd}}
\deleted{commutatif.}

% archive sheet number: 47

\scanpage{49}

% letterhead
{\footnotesize
\begin{tabular}{@{}l@{\quad}l@{}}
\deleted{MINISTERE DE L'EDUCATION NATIONALE} & \deleted{AVIGNON, le}\\
\deleted{ACADEMIE D'AIX-MARSEILLE} &\\
\hfil\rule{1.5cm}{0.4pt}\hfil &\\
\deleted{Inspection Académique} & \deleted{L'Inspecteur d'académie,}\\
\deleted{de Vaucluse} & \deleted{directeur des services départementaux}\\
\hfil\rule{1.5cm}{0.4pt}\hfil & \deleted{de l'Education Nationale de Vaucluse,}\\
\deleted{Tél. : 90 82 99 76} &\\
\deleted{Référence : Secrétariat n\textsuperscript{o}} & \deleted{à M}
\end{tabular}}
% end letterhead

\diagram{\begin{tikzcd}[ampersand replacement=\&,row sep=small]
{} \& \deleted{p\in\mathrm{Hom}(e,a)}\\
\deleted{\mathrm{Hom}(e,e)}\arrow[ur,"\deleted{\unsure{\widetilde p|\hat p|\varphi}}"'] \& {}
\end{tikzcd}}

\diagram{\begin{tikzcd}[ampersand replacement=\&,column sep=small,row sep=small,execute at end picture={\draw (current bounding box.south west)--(current bounding box.north);}]
{} \& \deleted{b}\arrow[dr,shift left=1.5,"\deleted{u}"]\arrow[dr,shift right=1.5,"\deleted{v}"'] \& {}\\
\deleted{e_a}\arrow[ur,shift left=1.5,"{\deleted{p_{b,u}}}"]\arrow[ur,shift right=1.5,"{\deleted{p_{b,v}}}"']\arrow[rr,"\deleted{p_a}"'] \& {} \& \deleted{a}
\end{tikzcd}}

\diagram{\begin{tikzcd}[ampersand replacement=\&,row sep=small]
\deleted{e}\arrow[r]\arrow[d,shift right=1.5] \& \deleted{a}\\
\deleted{e}\arrow[u,shift right=1.5]\arrow[ur,very thick] \& {}
\end{tikzcd}}

\[
\begin{aligned}
&\deleted{uq=p}\qquad\deleted{\text{\deleted{\unsure{avec}} avec } q\text{ comme }p}\\
&\hspace{5em}\deleted{\text{donc }uq=q\quad\text{donc }p=q}
\end{aligned}
\]

\[
\deleted{\unsure{q}}\qquad\deleted{\underbrace{uq}_{q}=\underbrace{vq'}_{q'}=p}
\]

\diagram{\begin{tikzcd}[ampersand replacement=\&,row sep=small,column sep=small]
{} \& \deleted{b}\arrow[dr,"\deleted{u}"] \& {}\\
\deleted{e}\arrow[ur,"\deleted{p_u}"]\arrow[dr,leftrightarrow,"\deleted{p_v}"'] \& {} \& \deleted{q}\\
{} \& \deleted{b}\arrow[ur,"\deleted{v}"] \& {}
\end{tikzcd}}

\[
\begin{aligned}
&\deleted{\deleted{\mathrm{Hom}(\unsure{x|b},a)}}\\
&\deleted{\mathrm{Hom}(x,a)}\\
&\deleted{x\longmapsto\mathrm{Hom}(e,x)}\qquad\deleted{p=qq'}\\
&\deleted{up=q}\qquad\deleted{q=v\varphi\deleted{\unsure{(w)}}\inserted{w'}}\\
&\deleted{p=v\bar v}\\
&\deleted{u'}
\end{aligned}
\]

\diagram{\begin{tikzcd}[ampersand replacement=\&,row sep=small,column sep=small]
{} \& {}\arrow[dr,dash,"\deleted{v}"] \& {}\\
{}\arrow[ur,dash]\arrow[rr,dash,"\deleted{p}"'] \& {} \& {}
\end{tikzcd}}

% archive sheet number: 48

\scanpage{50}
\mspage{\deleted{\unsure{5|4}}\unsure{4|5}}

\uline{Corollaire} Soit $X$ une catégorie finie, pseudo-cofiltr. \deleted{\unsure{et|est|sép.}}\inserted{(i.e $X^\wedge$ est loc. irréductible)}. Donc $\operatorname{Ord} X$ est un ensemble ordonné cofiltrant, admettant donc un plus petit élément. Soit $U$ l'ouvert de $X$ formé des objets de $X$, qui sont des plus petits objets. \deleted{On suppose} \inserted{Si} $U$ $1$-connexe, \deleted{alors} \deleted{\unsure{$X$|$U$}} est cofiltrant. (Donc si tout objet $0$-connexe de $X$ est $1$-connexe, alors $X$ est cofiltrant.)

En effet, il est immédiat \deleted{que} $U$ étant un ouvert cocofinal dans $X$, \inserted{(et $X$ étant pseudo cofiltrante)} que $X$ est cofiltrant dès que $U$ l'est. D'autre part, si $e\in U$, alors l'inclusion de la sous-catégorie pleine \inserted{$U_0$} de $X$ réduite à $e$ dans $U$ est une équivalence,\textsuperscript{*} ce qui implique que $U_0$ est $1$-connexe, $U$ l'étant, et que $U$ est cofiltrant sss $U_0$ l'est.
% slanted
\margin{Faux ! revoir la \uline{démonstration}}
\deleted{Donc} Il est immédiat aussi, $X$ étant pseudo-cofiltrant, que l'ouvert cocofinal $U$ l'est également, donc $U_0$ (équivalente à $U$) également.

Soit donc $M=\operatorname{End}_X(e)$. On peut appliquer à $M$ la proposition 1, qui implique que $M$ est cofiltrant, i.e. $U_0\simeq B_M$ est cofiltrant, qed

\uline{NB} Au lieu de l'hypothèse $X$ \uline{finie}, il suffisait de supposer que a) $X$ admet un plus petit objet $e$, et b) $\operatorname{End}_X(e)$ est fini.
% archive sheet number: 49

\scanpage{51}
\mspage{5}

\uline{Proposition 1} Soit $X$ une catégorie telle que

\deleted{\unsure{a)}} a) $X$ pseudocofiltrante \quad b) $X$ admet un plus petit objet $e$ \inserted{(conséquence de \unsure{a)} si $\unsure{\operatorname{Ord}|\operatorname{Ob}} X$ fini)} \deleted{objet minimal \unsure{$e$}} c) Pour tt $x\in\operatorname{Ob}X$, l'ens des sous-$\underbrace{\operatorname{Hom}(e,e)}_{M}$-ens. à droite de $\operatorname{Hom}(e,x)$ satisfait la condition minimale (ce qui est le cas si les $\operatorname{Hom}(e,x)$ sont finis). Alors

$1^\circ)$ $\forall x\in\operatorname{Ob}X$, parmi les sous-$M$-ens. à gauche de $\operatorname{Hom}(e,x)\overset{\mathrm{def}}=F(x)$ qui sont de la forme $fM$ (avec $f\in F(x)$), il y en a un plus petit, \inserted{soit $H_x=\bigcap_{f\in F(x)}fM$.} Il est formé de l'\unsure{ens|un} des $\lambda_{\deleted{\unsure{x}}}\in\operatorname{Hom}(e,x)$ \deleted{tels que} \inserted{qui sont} tels que \deleted{$\forall f\in\operatorname{Hom}($} $\forall f\in F(x)$, $\exists u\in M$ avec $\lambda=fu$.

$2^\circ)$ $H_x$ est donc un $M$-ens. à droite. De plus, pour tout $\deleted{\unsure{x}}f:x\longrightarrow y$ dans $X$, on a $H_y=f.H_x$. On a donc $\operatorname{card} H_y\leq\operatorname{card} H_x$ si $x\leq y$ ($\overset{\mathrm{def}}\Longleftrightarrow$ si $\operatorname{Hom}(x,y)\ne\varnothing$).

$3^\circ)$ Supposons les $F(x)$ finis (cf condition c) ci-dessus) et que $e$ soit aussi objet maximal, i.e. que $\operatorname{Ord} X$ soit réduit
% archive sheet number: 50

\scanpage{52}
\mspage{6}

à un élément. Alors pour tte flèche $f:x\longrightarrow y$ dans $X$, \deleted{$\lambda\unsure{\mapsto|-|\to} f$}
\[
\lambda\longmapsto f\lambda:\quad H_x\overset{\sim}{\longrightarrow} H_y
\]
est bijectif, (i.e les $H_x$ forment un syst. local sur $X$).

$4^\circ)$ Supposons, \deleted{\unsure{que}} en plus des conditions de a) b) c) et celles de $3^\circ)$, que $X$ soit $1$-connexe. Alors le système local des $H_x$ est \inserted{donc} constant, \deleted{\unsure{donc|dessus}} en particulier l'opération \inserted{à droite} de $M=\operatorname{Hom}(e,e)$ sur $H_x$ est constante, \deleted{i.e $\lambda\in$} et par suite l'
\[
\deleted{\lambda\in H_x,\quad u\in M\Longrightarrow\lambda u=\lambda}
\]
application
\[
u\longmapsto fu:\ H_e\longrightarrow H_x
\]
définie par un $f\in F(x):e\longrightarrow x$ ne dépend pas du choix de $f$, donc
\[
fu=gu\quad\text{pour }\begin{aligned}&u\in H_e,\\& f,g:e\rightrightarrows x.\end{aligned}
\]

% Printed form upside down relative to the mathematics.
% letterhead
{\footnotesize
\begin{tabular}{@{}l@{\quad}l@{}}
\deleted{MINISTERE DE L'EDUCATION NATIONALE} & \deleted{AVIGNON, le}\\
\deleted{ACADEMIE D'AIX-MARSEILLE} &\\
\hfil\rule{1.5cm}{0.4pt}\hfil &\\
\deleted{Inspection Académique} & \deleted{L'Inspecteur d'académie,}\\
\deleted{de Vaucluse} & \deleted{directeur des services départementaux}\\
\hfil\rule{1.5cm}{0.4pt}\hfil & \deleted{de l'Education Nationale de Vaucluse,}\\
\deleted{Tél. : 90 82 99 76} &\\
\deleted{Référence : Secrétariat n\textsuperscript{o}} & \deleted{à M}
\end{tabular}}
% end letterhead
% archive sheet number: 51
